\section{Trattazione matematica}\label{app:TrattazioneMatematica}
% Allegato: qui vanno i dettagli matematici che appesantirebbero il testo principale.
% Nel corpo della relazione rimandare qui con \ref{app:TrattazioneMatematica}
% (o con \cref{app:...} alle singole sottosezioni).
In questo allegato si riprendono in forma più rigorosa i modelli presentati nella \cref{sec:Prevedibilita}, con le definizioni, le proprietà principali e i calcoli degli esempi.

\subsection{Notazione}\label{app:Notazione}
\begin{itemize}
    \item \(\Sigma\): insieme dei caratteri ammessi (l'\emph{alfabeto}), con \(N = \vert\Sigma\vert\).
    \item \(L\): lunghezza della password.
    \item \(\Omega\): insieme di tutte le password possibili; nel modello di base \(\Omega = \Sigma^L\).
    \item \(X\): variabile aleatoria che rappresenta la password scelta dall'utente, a valori in \(\Omega\).
    \item \(P(x) = \mathbb{P}(X = x)\): probabilità che venga scelta la password \(x \in \Omega\).
    \item \(D\): variabile aleatoria che rappresenta l'informazione sull'utente in possesso dell'attaccante, a valori in un insieme \(\mathcal{D}\).
    \item \(\log_2\): logaritmo in base 2; tutte le entropie sono espresse in \emph{bit}.
    \item Si adotta la convenzione \(0 \log_2 0 = 0\).
\end{itemize}

\subsection{Modello teorico}\label{app:ModelloTeorico}
Il numero di password di lunghezza \(L\) su un alfabeto di \(N\) caratteri è
\[
    \vert\Omega\vert = \vert\Sigma^L\vert = N^L.
\]
Se si ammettono tutte le lunghezze da \(1\) a \(L\), il dominio diventa
\[
    \vert\Omega\vert = \sum_{\ell=1}^{L} N^\ell = \frac{N(N^L - 1)}{N - 1} \approx N^L \cdot \frac{N}{N-1},
\]
che per alfabeti realistici (\(N \geq 10\) caratteri ammessi) differisce da \(N^L\) per meno del \(12\%\). Questo significa che le password più corte contano poco, e \(N^L\) resta una buona approssimazione.

\paragraph{Entropia teorica.}
Se l'utente sceglie la password in modo uniforme, cioè \(P(x) = 1/\vert\Omega\vert\) per ogni \(x \in \Omega\), l'entropia è
\[
    \boxed{
        H = \log_2 \vert\Omega\vert = L \log_2 N.
    }
\]
Un attaccante che prova le password una alla volta ne prova al massimo \(2^H\), e in media la azzecca dopo \((2^H + 1)/2 \approx 2^{H-1}\). Ogni bit in più raddoppia quindi il lavoro dell'attaccante.

\paragraph{Tempo di attacco.}
Se l'attaccante prova \(v\) password al secondo, il tempo necessario nel caso peggiore è
\[
    T = \frac{2^H}{v}.
\]
È questa la grandezza riportata in \cref{fig:tempi-cracking-real}, con le ipotesi su \(v\) descritte nella \cref{sec:Prevedibilita}.
\begin{example}
    Con \(N = 94\) (caratteri stampabili) e \(L = 8\) si ha \(H = 8 \log_2 94 \approx 52{,}4\) bit, cioè circa \(6 \cdot 10^{15}\) password. Con \(v = 10^{10}\) tentativi al secondo servono \(T \approx 6 \cdot 10^{5}\) secondi, circa una settimana.
\end{example}
Nel caso peggiore servono sempre \(\vert\Omega\vert\) tentativi, qualunque sia la distribuzione. Il tempo \emph{atteso}, circa \(T/2\), è invece massimo quando la scelta è completamente casuale, perché nessuna distribuzione su \(\Omega\) ha entropia maggiore di quella uniforme: è quindi un limite massimo al tempo medio di attacco.

\subsection{Modello reale}\label{app:ModelloReale}
Le password reali non sono scelte in modo uniforme: i fattori di prevedibilità visti nella \cref{sec:Prevedibilita} concentrano la probabilità su una piccola parte di \(\Omega\). La distribuzione delle probabilità di \(X\) quindi non è più uniforme e va considerata generica.

\paragraph{Entropia di Shannon.}
La quantità \(-\log_2 P(x)\) è detta \emph{autoinformazione} della password \(x\): vale pochi bit se \(x\) è molto probabile, molti se è rara. L'entropia di \(X\) ne è il valore atteso:
\[
    \boxed{
        H(X) = -\sum_{x \in \Omega} P(x) \log_2 P(x) = \mathbb{E}\left[-\log_2 P(X)\right].
    }
\]
\begin{example}
    La difficoltà di una password non dipende dalla sua forma ma dalla sua probabilità. Una password \(x_0\) di 100 caratteri scelta dalla metà degli utenti ha \(P(x_0) = 1/2\), quindi autoinformazione \(-\log_2 \tfrac{1}{2} = 1\) bit: indovinarla è facile quanto indovinare il lancio di una moneta.
\end{example}

\paragraph{Limite superiore.}
Per ogni distribuzione su \(\Omega\)
\[
    H(X) \leq \log_2 \vert\Omega\vert = L \log_2 N.
\]

\paragraph{Entropia minima.}
L'entropia di Shannon è una media e non descrive bene l'attaccante che prova per prime le password più comuni. L'\emph{entropia minima} considera solo la password più probabile:
\[
    \boxed{
        H_{\infty}(X) = -\log_2 \left( \max_{x \in \Omega} P(x) \right).
    }
\]
La probabilità che l'attaccante indovini al primo tentativo, scegliendo la password più probabile, è esattamente \(2^{-H_{\infty}(X)}\). Inoltre valgono:
\[
    H_{\infty}(X) \leq H(X) \leq \log_2 \vert\Omega\vert.
\]

\begin{example}
    Si riprenda l'Esempio~\ref{ex:password-comune}: \(\vert\Omega\vert = 1000\), \(P(\texttt{12345678}) = 0{,}35\) e le restanti 999 password equiprobabili, ciascuna con probabilità \(0{,}65/999\). Allora
    \[
        \log_2 \vert\Omega\vert \approx 9{,}97 \text{ bit}, \qquad
        H(X) \approx 7{,}41 \text{ bit}, \qquad
        H_{\infty}(X) = -\log_2 0{,}35 \approx 1{,}51 \text{ bit}.
    \]
    L'entropia di Shannon resta relativamente alta, ma l'entropia minima mostra che l'attaccante indovina al primo tentativo con probabilità \(2^{-1{,}51} = 0{,}35\).
\end{example}

\subsection{Modello con dati sull'utente}\label{app:ModelloDati}
Sia \(D\) l'informazione sull'utente nota all'attaccante. In generale \(X\) e \(D\) non sono indipendenti: osservare \(D = d\) cambia la distribuzione della password da \(P(x)\) a \(P(x \mid d)\).

\paragraph{Entropia condizionata.}
L'incertezza residua su \(X\) dopo aver osservato \(D\) è
\[
    H(X \mid D) = \sum_{d \in \mathcal{D}} P(d) \, H(X \mid D = d),
    \qquad
    H(X \mid D = d) = -\sum_{x \in \Omega} P(x \mid d) \log_2 P(x \mid d).
\]

\paragraph{Informazione mutua.}
\[
    \boxed{
        I(X;D) = H(X) - H(X \mid D) = \sum_{x \in \Omega} \sum_{d \in \mathcal{D}} P(x, d) \log_2 \frac{P(x, d)}{P(x) P(d)}.
    }
\]
Si noti che \(I(X;D) \geq 0\), con uguaglianza se e solo se \(X\) e \(D\) sono indipendenti. Infatti l'ultima espressione è la divergenza di Kullback--Leibler tra \(P(x,d)\) e \(P(x)P(d)\), che è non negativa e nulla solo quando le due distribuzioni coincidono (disuguaglianza di Gibbs). Equivalentemente \(H(X \mid D) \leq H(X)\): \emph{in media} conoscere \(D\) non aumenta l'incertezza. Per un singolo valore \(d\) può invece accadere che \(H(X \mid D = d) > H(X)\); la disuguaglianza vale sulla media pesata da \(P(d)\).

\begin{example}
    \(X\) è uniforme su \(1024 = 2^{10}\) password, divise in 16 gruppi equiprobabili da 64; \(D\) indica il gruppo che contiene la password scelta. Dato \(D\), \(X\) è uniforme sulle 64 password del gruppo, quindi
    \[
        H(X) = 10 \text{ bit}, \qquad H(X \mid D) = \log_2 64 = 6 \text{ bit}, \qquad I(X;D) = 4 \text{ bit}.
    \]
    In questo caso \(I(X;D) = H(D) = \log_2 16\): poiché il gruppo è determinato dalla password, \(D\) rivela esattamente i 4 bit che servono a identificarlo.
\end{example}

\begin{example}
    Nell'Esempio~\ref{ex:telefono} i valori esatti sono \(H(X) = \log_2 10^9 \approx 29{,}90\) bit, \(H(X \mid D) = \log_2 10^3 \approx 9{,}97\) bit e \(I(X;D) = \log_2 10^6 \approx 19{,}93\) bit: gli arrotondamenti a \(30\), \(10\) e \(20\) bit sono corretti a meno di un decimo di bit, e l'informazione raccolta divide i tentativi per \(2^{I(X;D)} = 10^6\).
\end{example}

% Nelle sezioni precedenti si è visto come funzionano gli attacchi a dizionario (\cref{sec:AttacchiDizionario}) e come l'OSINT li renda più mirati (\cref{sec:IntegrazioneOSINT}).
% Resta da capire \emph{quanto} una password sia davvero difficile da indovinare.
% In questa sezione si introduce l'entropia, una misura dell'incertezza espressa in bit, e la si usa per confrontare tre scenari: password scelte a caso, password scelte da persone reali e password di cui l'attaccante conosce qualche informazione sull'utente.
% Vedremo che a ogni passaggio l'entropia diminuisce, e con essa il numero di tentativi necessari a un attacco.


% Modello di Markov e PCFG grammatiche probabilistiche

% \subsection{Passphrase}
% Si sa che gli umani tendenzialmente fanno schifo a inventarsi cose e che ricavano informazioni dall'ambiente etc..

% \[
%     P(w_1,\ldots,w_k)
%     =P(w_1)\prod_{i=2}^{k}P(w_i\mid w_1,\ldots,w_{i-1})
% \]
% \[
%     H(W_1,\ldots,W_k)=\sum_{i=1}^{k}H(W_i\mid W_1,\ldots,W_{i-1})
% \]

% [IDEA] Chiusura (3-4 righe, punto "impatto sugli attacchi"): l'esempio sopra in termini
%   pratici, 4 bit in meno = 16 volte meno tentativi. Messi insieme, schemi prevedibili e dati
%   OSINT trasformano una ricerca impossibile in un dizionario piccolo e mirato. Concludere
%   con il rimando alla sez. 6: le contromisure servono ad aumentare H(X|D), cioè a scegliere
%   password che non dipendono da ciò che si sa dell'utente (generatori, password manager).
